%\noindent 
In the previous parts we have examined agentic LLMs primarily through a technical lens: architectures, training procedures, reasoning strategies, tool use, and applications in various domains. In this part we shift gears. Rather than asking how agentic LLMs work and can be built, we ask how they \emph{behave}, socially. %Previous parts studied how agents act, this part studies how they act {\em together}. 
As agentic LLMs are increasingly deployed in social contexts (as assistants, collaborators, tutors, etc.), their behavioral inclinations start to matter practically. It has real-world consequences whether, for example, an agent tends to cooperate or defect in strategic situations, whether it responds to politeness, or whether it can track another agent's belief-states in a multi-agent simulation. 

Most research that we discuss in this part builds on methods originally developed to study humans and other animals, addressing strategic behavior, social etiquette, empathy, personality, and other traits. We end each chapter with a short reflection on the general question whether such methods can be productively applied to AI. 

In the next and final part, we go deeper into questions of LLM inner workings and the potential emergence of various forms of cognition, or even consciousness.
% The preceding four parts showed us how to create LLM Agents that reason, act, and interact with the world. This poses the question how such agents will behave in different contexts. Will they cooperate or compete, lead or follow? Is their behavior stable and predictable, or rather dynamic in novel ways? And what kind of behavior will emerge in larger systems of interacting LLM agents? The first chapter of this part will study different kinds of role-based and experimentally controlled behavior. The second chapter discusses emergent behavior in more open-ended settings. The final content chapter goes deeper, and will dive into parallels and differences between behaviour and cognition of LLM agents and humans. Finally, we will discuss a practical simulation of a society of Agents.
